\documentclass[11pt]{article}

\usepackage[a4paper,margin=1in]{geometry}
\usepackage{amsmath,amssymb}
\usepackage{graphicx}
\usepackage{booktabs}
\usepackage{multirow}
\usepackage{hyperref}
\usepackage{xcolor}
\usepackage{caption}
\usepackage{subcaption}
\usepackage{microtype}
% This paper uses Unicode characters from several scripts (Arabic-Indic,
% Devanagari, CJK, cuneiform). Compile with XeLaTeX or LuaLaTeX so the
% fontspec-based font selection below works; pdfLaTeX will not render
% these glyphs.
\usepackage{fontspec}
\usepackage{xeCJK}

% Inline font switches for non-Latin scripts. Each \newfontfamily picks a
% font that ships with macOS by default (Geeza Pro, Devanagari Sangam MN)
% plus Noto Sans Cuneiform for Babylonian wedges (install via
% `brew install --cask font-noto-sans-cuneiform`).
\newfontfamily\arabicscriptfont{Geeza Pro}[Scale=0.9]
\newfontfamily\devanagariscriptfont{Devanagari Sangam MN}[Scale=0.9]
\newfontfamily\cuneiformscriptfont{Noto Sans Cuneiform}[Scale=0.9]
\newcommand{\arabicnum}[1]{{\arabicscriptfont #1}}
\newcommand{\devnum}[1]{{\devanagariscriptfont #1}}
\newcommand{\cuneinum}[1]{{\cuneiformscriptfont #1}}

\hypersetup{
  colorlinks=true,
  linkcolor=blue!50!black,
  citecolor=blue!50!black,
  urlcolor=blue!60!black,
}

\title{Where the Number Helix Lives, and How to Find It:\\
       Cross-Architecture and Cross-Script Geometry of Integer Representations}

\author{Girish Gupta\thanks{\texttt{girish@girishgupta.com}}}
\date{\today}

\graphicspath{{out/}}

\begin{document}
\maketitle

\begin{abstract}
Kantamneni and Tegmark~\cite{kantamneni2025helix} showed that for an
integer $a \in [0, 99]$, the residual stream of three mid-sized LLMs
(GPT-J-6B, Pythia-6.9B, Llama-3.1-8B) encodes $a$ as a generalised
helix---a linear ``number-line'' axis plus four orthogonal modular
circles at periods $T \in \{2, 5, 10, 100\}$, and that this helix is
causally used to perform addition. Their basis itself is the LLM-scale
analogue of Nanda~et al.~\cite{nanda2023progress}, who reverse-engineered
the same trig solution in a small transformer trained from scratch on
modular addition---the canonical grokking task of
Power~et al.~\cite{power2022grokking}. We extend the LLM result along
three orthogonal axes. \textbf{Architecture}: we add a new model
family (Gemma~4, at E4B and 31B) and re-measure Pythia and Llama under
a wider protocol, recovering the helix on all four models but at
radically different depths---Pythia/Llama mid-stack, Gemma-4-E4B at
$L\!=\!5$ of $30$, Gemma-4-31B as a sharp spike at $L\!=\!29$ of $61$. \textbf{Script}: across eight numeral
systems we find a three-way split---positional base-10 scripts the
model has been trained on (Latin, Devanagari, Arabic-Indic, Persian,
positional Chinese) reproduce the helix on multilingual models;
additive systems (Roman, Greek alphabetic) collapse it in two
qualitatively distinct ways distinguishable only by inspecting where
PC1 and helix $R^2$ peak in layer-space; and Babylonian cuneiform, a
positional base-60 system, encodes its native period~60 on
\emph{all four models}, provided we widen the input window past one
base-60 wrap and add $T=60$ to the trig basis (adding the single
missing period recovers $7\!-\!13$ percentage points of variance
uniformly across the four models). 2-D PCA scatters of the residual
stream further reveal that on Gemma the first principal component is
not magnitude but \emph{tokenisation} (PC1 captures up to $96.6\%$ of
variance but separates points by digit count, not value).
We extract two methodological lessons: (M1) any basis-dependent metric
must be cross-checked against a basis-free FFT and against a 2-D PCA
scatter; (M2) the input range must contain several wraps of every
candidate period.
\end{abstract}

\section{Introduction}
\label{sec:introduction}

Mechanistic interpretability has produced a growing catalogue of
``cognitive'' structures embedded in transformer residual
streams---induction heads, indirect-object identification circuits,
and so on. Arithmetic has been a particularly load-bearing case:
Power~et al.~\cite{power2022grokking} observed that a small transformer
trained on modular addition continues training long past memorisation
before suddenly generalising (``grokking''), and Nanda~et
al.~\cite{nanda2023progress} reverse-engineered the resulting solution
as a Fourier/trig basis---integers laid out on circles, addition
implemented via $\cos(a+b) = \cos a \cos b - \sin a \sin b$. The
number-helix of Kantamneni and Tegmark \cite{kantamneni2025helix}
carries the same basis into \emph{pretrained} LLMs on the naturals:
integers are not stored as opaque embeddings but laid out as a 3-D
spiral---a linear number-line axis plus four orthogonal circles---and
addition is performed by rotating the spiral's circular components.
The claim is striking; it is also based on three models in the same
broad decoder-only lineage (GPT-J-6B, Pythia-6.9B, Llama-3.1-8B) and
on a single numeral system (Latin digits).

This paper asks: \emph{how general is the helix?} We ask the question
along two independent axes:

\begin{enumerate}
  \item \textbf{Architecture}: does the helix extend beyond the
    decoder-only lineage Kantamneni and Tegmark studied? We add
    Gemma~4 (E4B and 31B---different tokeniser, different training
    mix, different depth) and re-measure Pythia and Llama under a
    wider protocol. If the helix is there, at what depth?
  \item \textbf{Script}: when the model reads ``\arabicnum{٢٣}'',
    ``\devnum{२३}'', ``\cuneinum{𒌋𒌋𒁹𒁹𒁹}'', does it represent
    these as numbers in the same helical way, or as ordinary
    multi-byte words?
\end{enumerate}

Both axes turn out to matter, and both surface methodological
subtleties that are easy to miss if one runs only the paper's exact
measurement recipe. Section~\ref{sec:methods} reviews the helix
construction and our two basis-free diagnostics. Section~\ref{sec:arch}
reports the layer-depth findings across four models.
Section~\ref{sec:script} compares eight numeral systems, and analyses
the two qualitatively distinct ways the helix collapses on
non-positional scripts. Section~\ref{sec:babylonian} contains the
main scientific surprise: Babylonian cuneiform is represented at its
native base 60 on every model we tested, once the measurement window
is widened past a single wrap and the basis is extended with $T=60$.
Section~\ref{sec:method-matters} distills the two
generalisable methodological lessons.

\section{Methods}
\label{sec:methods}

\subsection{Setup}
For each integer $a \in [0, n_{\max})$ we feed the string
``\verb*| <numeral(a)>|'' to the model and pool the residual stream at
a chosen layer $L$ across the sub-tokens of the numeral. With $d$ the
hidden width, the result is an activation matrix
$H \in \mathbb{R}^{n_{\max} \times d}$. Latin digits 0--99 are single
tokens on Pythia, so pooling is trivial there; in every other case
(non-Latin scripts; Latin on Llama and Gemma; the $n_{\max}\!=\!600$
sweep) we average activations across the numeral's sub-tokens to avoid
a last-byte cycling artifact in which $\{23, 33, 43, \dots\}$ read
identical activations because they end in the same glyph.

\subsection{The helix basis}
The Kantamneni--Tegmark hypothesis is that $h(a)$, for fixed $L$,
factors as
\begin{equation}
\label{eq:helix}
h(a) \;\approx\; W^{\top}\,B(a), \qquad
B(a) \;=\; \bigl[a,\;\cos\tfrac{2\pi a}{T_1},\,\sin\tfrac{2\pi a}{T_1},\,
\dots,\,\cos\tfrac{2\pi a}{T_K},\,\sin\tfrac{2\pi a}{T_K}\bigr],
\end{equation}
with $T = \{2, 5, 10, 100\}$ for Latin in $[0,99]$, giving a
$(1\!+\!2K) = 9$-dimensional feature space. The linear coordinate is
the ``rise'' of the helix; each $(\cos, \sin)$ pair places $a$ on the
unit circle at angle $2\pi a / T$, so integers sharing the same
$a \bmod T$ land on the same point. Addition becomes rotation of
angles: the paper's ``Clock'' algorithm.

We solve $H \approx B W$ by ordinary least squares and report the
variance-weighted coefficient of determination, which we call
\emph{helix $R^2$}.

\subsection{Three diagnostics}
\label{sec:diagnostics}
Helix $R^2$ alone is fragile---it depends on the fixed basis $T$. We
pair it with two basis-free diagnostics.

\paragraph{Fourier diagnostic.}
For each hidden dimension $j$, the column $H_{\cdot,j}$ is a sequence
indexed by $a$. We take its discrete Fourier transform, average the
magnitudes across $j$, and inspect the spectrum. Periodic structure
\emph{at any period}, including ones not in $T$, surfaces as a peak.
We label the top peaks via \texttt{scipy.signal.find\_peaks} so that
the FFT panel honestly reports whatever periods the model actually
encodes, rather than only those the paper's basis assumes.

\paragraph{PC1 diagnostic.}
The leading principal component of $H$ is the direction of maximum
variance. If the helix's linear axis is present, plotting PC1 against
$a$ should yield a clean line; we report the linear-fit $R^2$ as
\emph{PC1 $R^2$}.

\paragraph{Eckart--Young upper bound.}
Helix $R^2$ uses a $9$-D fit, but $9$ is also the dimensionality of
many other regression targets. The best possible $9$-D fit is the
$9$-D PCA reconstruction; we call its variance-weighted $R^2$ the
\emph{$K$-D PCA bound}, and report the ratio $\rho = R^2_{\text{helix}}
/ R^2_{\text{PCA}}$ as \emph{helix/PCA}. When $\rho \gtrsim 0.8$, the
trig basis is \emph{the} $9$-D structure of the residual stream rather
than merely \emph{a} $9$-D approximation of it.

\subsection{Layer sweep}
The paper reads at layer $L = N/2$ in an $N$-layer model. We capture
activations at all $L = 0, 1, \dots, N$ in one forward pass per
integer and compute (PC1 $R^2$, helix $R^2$, $K$-D PCA bound) at each.
The peak helix-$R^2$ layer is then used for the standard regression
fit. As Section~\ref{sec:arch} shows, ``read the middle layer'' is a
Pythia/GPT-J convention, not a universal one.

\subsection{2-D PCA scatter}
\label{sec:pca2d-method}
At the helix-peak layer of each cell we scatter the first two
principal components of $H$, colour-coded by $a$, with consecutive
integers connected by a faint trajectory line. This is the cheapest
basis-free diagnostic for \emph{non}-periodic structure (Roman's
threshold-aligned clusters; Greek's diffuse noise; Gemma's
tokenisation-dominated PC1; Section~\ref{sec:pca2d}).

\section{Architecture: the helix is consistent, the layer is not}
\label{sec:arch}

We test four base (non-instruction-tuned) models:
Pythia-6.9B, Llama-3.1-8B, Gemma-4-E4B (efficient $\sim$4\,B), and
Gemma-4-31B. All four converge on the same geometric encoding on
Latin digits 0--99 (Table~\ref{tab:arch}), but the depth at which the
helix peaks varies dramatically, and so does the shape of the
depth profile (Figure~\ref{fig:layer-sweep}).

\begin{table}[h!]
\centering
\caption{Latin helix metrics at the peak layer of each model.
``helix/PCA'' is the ratio of helix $R^2$ to the $9$-D PCA upper bound
(values $\gtrsim\!0.8$ indicate the trig basis dominates the
$9$-D subspace).}
\label{tab:arch}
\begin{tabular}{lcccc}
\toprule
Model & helix $R^2$ & helix/PCA & peak layer & total layers \\
\midrule
Pythia-6.9B   & 0.53 & 0.85 & 15 & 32 \\
Llama-3.1-8B  & 0.63 & 0.85 & 14 & 32 \\
Gemma-4-E4B   & 0.50 & 0.77 & 5  & $\sim$30 \\
Gemma-4-31B   & \textbf{0.70} & 0.75 & 29 & 61 \\
\bottomrule
\end{tabular}
\end{table}

\begin{figure}[h!]
\centering
\includegraphics[width=0.95\linewidth]{_compare/fig_latin_sweep_4models.png}
\caption{PC1 $R^2$, helix $R^2$ and helix/PCA ratio plotted against
layer index, for each of the four models (Latin digits 0--99). Pythia
and Llama rise into the back half of the stack and stay high.
Gemma-4-E4B peaks early and decays. Gemma-4-31B is nearly flat at
$\sim$0.5 with a sharp spike at $L\!=\!29$ that reads as ``no helix''
if one samples at $L \!=\! 30 = N/2$.}
\label{fig:layer-sweep}
\end{figure}

Three depth regimes emerge:
\begin{itemize}
  \item \textbf{Pythia and Llama}: helix forms in the back half of the
    stack and remains until at least $L = N-1$. Either $L = N/2$ or
    $L = N-1$ would have worked.
  \item \textbf{Gemma-4-E4B}: helix is strongest at $L = 0\!-\!5$ and
    decays in later layers. The paper's middle-layer recipe would
    literally read at decayed structure.
  \item \textbf{Gemma-4-31B}: helix $R^2$ is essentially flat at
    $\sim$0.5 across the whole stack, with a single sharp peak of 0.70
    at $L = 29$. The default $L = N/2 = 30$ misses the peak by exactly
    one layer.
\end{itemize}
We initially concluded ``Gemma-4-31B does not form a helix'' from a
default-layer reading---a conclusion wrong by a single layer. The
methodological consequence is unambiguous: \emph{for any new model, a
layer sweep is mandatory.} The helix is general; the recipe for
finding it is not.

\section{Script: positional structure is necessary, training exposure
is what makes the helix appear}
\label{sec:script}

\subsection{The numeral systems at a glance}
For readers unfamiliar with the non-Latin systems we test:
\begin{itemize}
  \item \textbf{Arabic-Indic} (\arabicnum{٠١٢٣٤٥٦٧٨٩}) and
    \textbf{Persian} (\arabicnum{۰۱۲۳۴۵۶۷۸۹}) are positional
    base-10 like Latin, differing only in glyph shape (Persian is a
    visual variant of Arabic-Indic). Numerals are written
    most-significant-first despite Arabic script flowing
    right-to-left.
  \item \textbf{Devanagari} (\devnum{०१२३४५६७८९}) is also positional
    base-10, used across Indian languages including Hindi, Sanskrit,
    and Marathi.
  \item \textbf{Positional Chinese} uses the digits
    \mbox{零一二三四五六七八九} in the same place-value layout as
    Latin (e.g.\ \mbox{二三} for 23). Classical Chinese also admits a
    non-positional form (\mbox{二十三}, ``two-tens-three''); we do
    not use it here.
  \item \textbf{Greek alphabetic (Milesian)} is additive: each letter
    carries a fixed value ($\alpha = 1, \dots, \theta = 9,
    \iota = 10, \kappa = 20, \dots, \pi = 80, \varkappa = 90,
    \rho = 100, \dots$). 23 is $\kappa\gamma = 20 + 3$. There is no
    shared ``ones digit'' between 13, 23, and 33.
  \item \textbf{Roman} is additive with subtractive shortcuts at 4,
    9, 40, 90, etc.: I$=1$, V$=5$, X$=10$, L$=50$, C$=100$.
    23 = XXIII. Glyphs shift at $4\leftrightarrow 5$,
    $9\leftrightarrow 10$, $49\leftrightarrow 50$,
    $89\leftrightarrow 90$ --- the same boundaries that later appear
    as staircase discontinuities in the PC1 panel
    (Section~\ref{sec:script-collapse}).
  \item \textbf{Babylonian cuneiform} is positional at base 60,
    additive within each sexagesimal column. Two glyphs combine
    inside a column: \cuneinum{𒁹} (one) and \cuneinum{𒌋} (ten).
    23 is \cuneinum{𒌋𒌋𒁹𒁹𒁹} --- two tens plus three ones, all
    in one column. 60 opens a new column to the left
    (Section~\ref{sec:babylonian}).
\end{itemize}

\subsection{Setup}
We rendered every integer in $[0,99]$ in eight numeral systems
(Table~\ref{tab:scripts}) and re-ran the analysis under mean-pooling.

\begin{table}[h!]
\centering
\caption{Numeral systems tested. The Babylonian column is the most
interesting and is treated separately in Section~\ref{sec:babylonian}.}
\label{tab:scripts}
\small
\begin{tabular}{lll}
\toprule
Script & Example (23) & Family \\
\midrule
Latin             & \texttt{23}              & positional base-10 \\
Arabic-Indic      & \arabicnum{٢٣}              & positional base-10 \\
Persian           & \arabicnum{۲۳}              & positional base-10 \\
Devanagari        & \devnum{२३}          & positional base-10 \\
Chinese (pos.)    & 二三                      & positional base-10 \\
Greek alphabetic  & \(\kappa\gamma\)         & additive (Milesian) \\
Roman             & \texttt{XXIII}           & additive \\
Babylonian        & \cuneinum{𒌋𒌋𒁹𒁹𒁹}     & positional base-60, additive within column \\
\bottomrule
\end{tabular}
\end{table}

\subsection{Helix/PCA across scripts}
\begin{table}[h!]
\centering
\caption{Helix/PCA ratio at each model's peak layer, by script (mean
pooling). Bold = highest score per model.}
\label{tab:script-results}
\begin{tabular}{lcccc}
\toprule
Script & Pythia-6.9B & Llama-3.1-8B & Gemma-4-E4B & Gemma-4-31B \\
\midrule
Latin             & \textbf{0.85} & \textbf{0.85} & \textbf{0.77} & 0.75 \\
Arabic-Indic      & 0.55 & 0.58 & 0.66 & 0.69 \\
Persian           & 0.55 & 0.70 & 0.72 & 0.72 \\
Devanagari        & 0.53 & 0.61 & 0.78 & \textbf{0.79} \\
Chinese           & 0.60 & 0.58 & 0.68 & 0.66 \\
Greek (additive)  & 0.54 & 0.57 & 0.57 & 0.57 \\
Roman  (additive) & 0.54 & 0.55 & 0.62 & 0.64 \\
\bottomrule
\end{tabular}
\end{table}

Three observations stand out:

\paragraph{Non-positional scripts collapse uniformly.} On every model,
Greek and Roman collapse the helix/PCA ratio to $\sim$0.55--0.65.
A trig basis with a tens-loop only buys representational efficiency if
the numeral system has a tens column. Roman \texttt{IV} versus
\texttt{V} is not the kind of transition a fixed-period circle can
align to.

\paragraph{Latin is uniform.} All four models clear 0.75 on Latin. The
shape itself is not model-specific.

\paragraph{Non-Latin positional cells split by training data.} Pythia
on Devanagari (0.53) is statistically indistinguishable from Pythia on
Greek (0.54): the helix is gone. Llama on Devanagari (0.61) is a
little better. Gemma-4-31B on Devanagari hits 0.79, matching its own
Latin score of 0.75. Pythia's training mix is English-heavy; Llama
saw more multilingual data; Gemma is the most multilingual of the
four. \emph{Positional structure is necessary for the helix, but
training exposure to the specific script is what makes it appear.}
The encoding is built per-script during pre-training; architecture
alone is not enough.

\subsection{Two qualitatively distinct collapses on non-positional
scripts}
\label{sec:script-collapse}
The naive reading of Table~\ref{tab:script-results} is that ``Greek
and Roman collapse the helix.'' That is correct on the numerical
metric but misleading: the underlying geometry is qualitatively
different in the two cases. Reading PC1 $R^2$ alongside helix $R^2$,
\emph{and the layer at which each peaks}, reveals the difference
(Table~\ref{tab:script-collapse}).

\begin{table}[h!]
\centering
\caption{Layer of peak PC1 and helix $R^2$ on Latin, Roman, Greek.
Latin: both peak at the same layer (a single helix). Roman:
peaks are 9--20 layers apart (decomposed parallel modules). Greek:
helix $R^2$ peaks at $L=0$ (tokeniser artifact at the embedding).}
\label{tab:script-collapse}
\begin{tabular}{llccc}
\toprule
Model & Script & PC1 $R^2$ peak & helix $R^2$ peak & layer gap \\
\midrule
Pythia-6.9B  & Latin & 0.96 @ L5  & 0.53 @ L32 & --- \\
Pythia-6.9B  & Roman & 0.66 @ L9  & 0.36 @ L18 & \textbf{9} \\
Pythia-6.9B  & Greek & 0.10 @ L9  & 0.28 @ \textbf{L0}  & (artifact) \\
Llama-3.1-8B & Roman & 0.51 @ L16 & 0.36 @ L17 & 1 \\
Llama-3.1-8B & Greek & 0.17 @ L5  & 0.27 @ \textbf{L0}  & (artifact) \\
Gemma-4-E4B  & Roman & 0.54 @ L4  & 0.61 @ L18 & \textbf{14} \\
Gemma-4-31B  & Roman & 0.49 @ L11 & 0.64 @ L31 & \textbf{20} \\
\bottomrule
\end{tabular}
\end{table}

\paragraph{Roman: a decomposed encoding, not a fused helix.}
On Latin, PC1 and helix $R^2$ peak at the same mid-stack layer;
magnitude and periodicity co-locate, and the trig basis literally fits
a helix (line $\times$ circles). On Roman, PC1 and helix $R^2$ peak at
\emph{different} layers---9 to 20 apart. Magnitude is present
(PC1 $R^2$ climbs to 0.49--0.66, so the model does coarsely order
Roman numerals), and so is periodic structure (Figure~\ref{fig:roman}
shows clean FFT peaks at $T=10$ for the X-cycle and $T=5$ for the
V-cycle), but they live in different parts of the residual stream.
Geometrically, Roman is encoded as a \emph{number-line module} at one
depth plus a \emph{letter-cycling module} at another, not as a fused
helix. The PC1-vs-$a$ panel of Figure~\ref{fig:roman} is also
informative: it is not a straight line but a staircase with
discontinuities at the I/V, IX/X, XL/L, and XC/C boundaries---the
underlying signal is piecewise linear, not linear, which is why the
naive linear-fit $R^2 = 0.23$ understates the magnitude information
that is present.

\begin{figure}[h!]
\centering
\includegraphics[width=0.85\linewidth]{EleutherAI__pythia-6.9b/roman/mean/fig2_fourier_pc1.png}
\caption{Pythia-6.9B on Roman numerals. \textbf{Top}: FFT spectrum
shows clean peaks at $T = 10$ (X cycle) and $T = 5$ (V cycle)---the
visual periodicity of the additive system. \textbf{Bottom}: PC1 versus
$a$ is a piecewise function with jumps at the Roman threshold values,
not a line. The two components live in different layers
(Table~\ref{tab:script-collapse}).}
\label{fig:roman}
\end{figure}

\paragraph{Greek: no quantitative representation at all.}
PC1 $R^2$ peaks at $0.10$--$0.21$ on every model; the linear axis
never appears. Helix $R^2$ peaks at layer $0$ on three of four
models---i.e.\ at the output of the token embedding, before any
transformer block has executed (Figure~\ref{fig:greek}). The FFT
magnitudes are $\sim$0.20 (compare $\sim$30 for Latin); the apparent
$T=10$ peak is a tokeniser artifact in which the second Greek letter
cycles through $\alpha, \beta, \gamma, \dots, \theta$ every ten
numbers. There is no number representation at any layer for the
``helix collapse'' to collapse \emph{from}.

\begin{figure}[h!]
\centering
\includegraphics[width=0.85\linewidth]{EleutherAI__pythia-6.9b/greek/mean/fig2_fourier_pc1.png}
\caption{Pythia-6.9B on Greek alphabetic numerals. The FFT magnitude
is two orders of magnitude smaller than on Latin; the only periodic
structure is a $T=10$ artifact at $L=0$ from the alphabetic
construction $\kappa + \gamma = $ \(\kappa\gamma\). PC1 $R^2 = 0.004$:
the model is not representing Greek numerals as quantities.}
\label{fig:greek}
\end{figure}

The two non-positional collapses are not the same failure mode. Roman
is encoded as a sum of separately-stored parts; Greek is not encoded
quantitatively at all. The standard ``helix $R^2$'' metric averages
over both modes and reports them as comparable; reading PC1 $R^2$ and
its layer-of-peak alongside helix $R^2$ is what separates them.

\subsection{2D PCA: cluster structure made visible (and the PC1-is-tokenization surprise)}
\label{sec:pca2d}

To probe non-periodic structure that the FFT and the trig-basis fit
cannot see, we added a basis-free diagnostic to every cell: a 2D PCA
scatter of $h(a)$ at the helix-peak layer, coloured by $a$, with each
numeral labelled and consecutive integers connected by a faint
trajectory line. Eight scripts $\times$ four models $=$ 32 panels at
\texttt{out/<model>/<script>/mean/fig4\_pca\_2d.png}. Three findings
fall out.

\textbf{Roman: cluster-by-prefix is visible directly.} The 2D PCA
for Pythia/Roman (Figure~\ref{fig:roman_pca}) shows distinct regions
of representation space corresponding to the Roman-numeral prefix
classes (I/V, X, L, XC), with the trajectory line zigzagging between
clusters at the threshold integers ($4 \to 5$, $9 \to 10$,
$49 \to 50$, $89 \to 90$). Roman's representation is not a manifold
but a partitioning of representation space into prefix classes with
crude ordering within each class.

\begin{figure}[h!]
\centering
\includegraphics[width=0.7\linewidth]{EleutherAI__pythia-6.9b/roman/mean/fig4_pca_2d.png}
\caption{Pythia-6.9B on Roman numerals. 2D PCA at the helix-peak
layer reveals discrete clusters by Roman-numeral prefix (I/V, X, L,
XC), with the trajectory line zigzagging at threshold integers. The
representation is a partition, not a continuous manifold.}
\label{fig:roman_pca}
\end{figure}

\textbf{Greek: diffuse noise, confirming the no-representation
reading.} PC1 + PC2 $= 15.5\%$ of total variance on Pythia/Greek (less
than half of any other cell) and the points scatter without
recognisable shape. There is no learned representation to find; the
``helix collapse'' was a collapse of nothing.

\textbf{The Gemma surprise: PC1 is not magnitude; it is token shape.}
On Gemma-4-31B, helix/PCA $= 0.75$ and helix $R^2 = 0.70$ on Latin
---strong by any measure. But PC1 captures \textbf{70.9\%} of variance
(versus Pythia's 20.4\%) and that variance is along token-shape
rather than magnitude (Figure~\ref{fig:gemma_latin_pca}). Single-digit
numerals (0--9) form one cluster, two-digit numerals starting with
``1'' (10--19) form another, and 20--99 fill a third dense region.
The split is by digit count and leading-digit class. The helix lives
in lower-variance directions; PC1 alone would not have found it.

\begin{figure}[h!]
\centering
\includegraphics[width=0.7\linewidth]{google__gemma-4-31B/latin/mean/fig4_pca_2d.png}
\caption{Gemma-4-31B on Latin digits. PC1 captures 70.9\% of variance
but separates the points by \emph{token shape} (single-digit
vs.~teens vs.~20--99), not magnitude. Helix $R^2 = 0.70$ persists at
this layer, in lower-variance directions PC1 does not reach.}
\label{fig:gemma_latin_pca}
\end{figure}

The same pattern appears on Devanagari with PC1 $= 94.6\%$ (split
between $०$--$९$ and $१०$--$९९$), and on Roman at PC1 $= 96.6\%$
(split at the L threshold). In each case PC1 captures massive
variance, but the variance is overwhelmingly tokenization features,
not the number line. In Pythia, by contrast, PC1 $R^2$ and
linearity-in-$a$ correspond cleanly: looking at PC1 is looking at
magnitude.

This carries one specific methodological warning: \textbf{a high PC1
$R^2$ is not by itself evidence of a clean magnitude axis}---one must
also inspect what PC1 separates. In larger or more multilingual
models, tokenization features may dominate the first principal
component while the helix occupies smaller-variance subspaces. This
also partly explains the depth-variance finding of
Section~\ref{sec:arch}: Gemma's helix shares residual-stream
space with strong tokenization signals, so it appears as a sharper
spike at a specific layer rather than a broad mid-stack plateau.

\section{Babylonian: base 60 is there, once you measure properly}
\label{sec:babylonian}

Babylonian cuneiform is positional at base 60, additive within each
sexagesimal column. The number 23 is \cuneinum{𒌋𒌋𒁹𒁹𒁹} (two tens
plus three ones); 60 is \cuneinum{𒁹\,𒑊} (one in the 60s column, a
late-period zero placeholder in the ones column); 599 is
\cuneinum{𒁹𒁹𒁹𒁹𒁹𒁹𒁹𒁹𒁹\,𒌋𒌋𒌋𒌋𒌋𒁹𒁹𒁹𒁹𒁹𒁹𒁹𒁹𒁹}. The system has all the
ingredients of positional notation, at base 60 instead of base 10.

\subsection{The narrow-window measurement says ``no''---but cannot}
Running the standard $n_{\max} = 100$ protocol with the paper's basis
$T = \{2, 5, 10, 100\}$ gives Table~\ref{tab:bab-n100}.

\begin{table}[h!]
\centering
\caption{Babylonian on the standard 0--99 window, paper basis.}
\label{tab:bab-n100}
\begin{tabular}{lccc}
\toprule
Model & PC1 $R^2$ & helix $R^2$ & helix/PCA \\
\midrule
Pythia-6.9B   & 0.06 & 0.71 & 0.74 \\
Llama-3.1-8B  & 0.10 & 0.65 & 0.69 \\
Gemma-4-E4B   & 0.40 & 0.75 & 0.77 \\
Gemma-4-31B   & 0.51 & 0.76 & 0.77 \\
\bottomrule
\end{tabular}
\end{table}

A naive read is that Pythia and Llama have circles but no number-line
(PC1 $R^2 \approx 0$), while Gemma has both. Inspecting the FFT
(Figure~\ref{fig:bab-n100}) shows labelled peaks at $T = 10$ and
$T = 100$, but \emph{no} $T = 60$. The PC1 panel shows two clusters
with a hard step at 60---the wrap-boundary is encoded as a categorical
state change, not as a rotation.

\begin{figure}[h!]
\centering
\includegraphics[width=0.85\linewidth]{google__gemma-4-31B/babylonian/mean/fig2_fourier_pc1.png}
\caption{Gemma-4-31B on Babylonian at $n_{\max}=100$. Top FFT shows
$T=10$ and $T=100$; \emph{no} $T=60$. PC1 (bottom) shows numbers
0--59 versus 60--99 as two clusters with a step between them.}
\label{fig:bab-n100}
\end{figure}

This is the measurement artifact: in $[0,99]$, base 60 wraps exactly
\emph{once}, at $a = 60$. A single discontinuity is not a periodic
signal an FFT can localise. ``No base-60 structure'' is, at this
$n_{\max}$, indistinguishable from ``structure the window cannot
resolve.''

\subsection{Widening the window: $T \approx 60$ emerges in the FFT}

At $n_{\max} = 600$---ten full sexagesimal wraps---the picture changes
sharply (Table~\ref{tab:bab-n600}). Every model now has a real
magnitude axis (PC1 $R^2 \approx 0.55$--$0.69$); Pythia's and Llama's
previous PC1 $R^2 \approx 0$ was an artifact of one wrap.

\begin{table}[h!]
\centering
\caption{Babylonian, $n_{\max}=100 \to 600$ (paper basis $T = \{2,5,10,100\}$).
Both PC1 and FFT diagnostics improve; helix $R^2$ \emph{drops} because
$T = 60$ is not in the basis.}
\label{tab:bab-n600}
\begin{tabular}{lccc}
\toprule
Model & PC1 $R^2$ ($100\!\to\!600$) & helix $R^2$ ($100\!\to\!600$) & helix/PCA \\
\midrule
Pythia-6.9B   & 0.06 $\to$ \textbf{0.55} & 0.71 $\to$ 0.66 & 0.74 $\to$ 0.69 \\
Llama-3.1-8B  & 0.10 $\to$ \textbf{0.55} & 0.65 $\to$ 0.52 & 0.69 $\to$ 0.56 \\
Gemma-4-E4B   & 0.40 $\to$ \textbf{0.69} & 0.75 $\to$ 0.61 & 0.77 $\to$ 0.64 \\
Gemma-4-31B   & 0.51 $\to$ \textbf{0.64} & 0.76 $\to$ 0.50 & 0.77 $\to$ 0.52 \\
\bottomrule
\end{tabular}
\end{table}

The FFT at $n_{\max}=600$ shows a clean labelled peak at $T \approx 60$
(Figure~\ref{fig:bab-n600-pythia}). For Llama and Gemma-4-31B, base-60
structure appears mostly via its second harmonic at $T = 30$
(frequency $2/60 = 1/30$, Figure~\ref{fig:bab-n600-gemma}). A pure
sinusoid at $T=60$ would have only the fundamental; the presence of
strong even harmonics indicates a more sawtooth-like underlying
signal, consistent with the per-column reset structure of cuneiform.

\begin{figure}[h!]
\centering
\includegraphics[width=0.85\linewidth]{EleutherAI__pythia-6.9b/babylonian/mean_n600/fig2_fourier_pc1.png}
\caption{Pythia-6.9B on Babylonian at $n_{\max}=600$. A clean labelled
peak at $T \approx 60$ appears alongside $T = 10$ and various
harmonics. The PC1 panel shows the magnitude axis with a per-column
sawtooth---the model resets into a new sexagesimal column every 60
integers.}
\label{fig:bab-n600-pythia}
\end{figure}

\begin{figure}[h!]
\centering
\includegraphics[width=0.85\linewidth]{google__gemma-4-31B/babylonian/mean_n600/fig2_fourier_pc1.png}
\caption{Gemma-4-31B on Babylonian at $n_{\max}=600$. Base-60 structure
manifests primarily through its second harmonic at $T=30$, indicating
a sawtooth-like signal rather than a pure sinusoid.}
\label{fig:bab-n600-gemma}
\end{figure}

So the structure is \emph{present in the FFT}---but helix $R^2$ on the
paper basis \emph{drops} (Table~\ref{tab:bab-n600}), because
$T \in \{2, 5, 10, 100\}$ has no column against which to regress
period-60 variance. The basis is blind to a frequency the data
contains.

\subsection{An extended basis recovers $7$--$13$ percentage points of
variance, on every model}
\label{sec:bab-extended}
We re-ran the $n_{\max}=600$ sweep with $T = \{2, 5, 10, 60, 100\}$---a
single additional period, two additional columns
($\cos\frac{2\pi a}{60}, \sin\frac{2\pi a}{60}$).
Table~\ref{tab:bab-extended} reports the result.

\begin{table}[h!]
\centering
\caption{Babylonian at $n_{\max}=600$: paper basis $T=\{2,5,10,100\}$
vs.\ extended basis $T=\{2,5,10,60,100\}$. Adding the single missing
period recovers $7$--$13$ percentage points of variance on every model.}
\label{tab:bab-extended}
\begin{tabular}{lcc cc}
\toprule
Model & helix $R^2$ (paper $\to$ ext.) & $\Delta$ & helix/PCA (paper $\to$ ext.) & $\Delta$ \\
\midrule
Pythia-6.9B   & 0.661 $\to$ \textbf{0.747} & \textbf{+0.086} & 0.687 $\to$ 0.772 & +0.085 \\
Llama-3.1-8B  & 0.515 $\to$ \textbf{0.622} & \textbf{+0.107} & 0.555 $\to$ 0.645 & +0.090 \\
Gemma-4-E4B   & 0.611 $\to$ \textbf{0.698} & \textbf{+0.087} & 0.637 $\to$ 0.699 & +0.062 \\
Gemma-4-31B   & 0.501 $\to$ \textbf{0.634} & \textbf{+0.133} & 0.523 $\to$ 0.636 & +0.113 \\
\bottomrule
\end{tabular}
\end{table}

The two models whose FFT showed strongest second-harmonic content
(Llama-3.1-8B and Gemma-4-31B) gain the most ($+0.107$ and $+0.133$).
The regression cannot extract a non-sinusoidal signal from a basis
that lacks its fundamental; providing $T=60$ anchors the phase and
the model's sawtooth-like representation lands on it. PC1 $R^2$ is
unchanged by the basis change, confirming that no other variable has
shifted.

The conclusion is unambiguous:
\begin{quote}
\textbf{Base-60 structure is present in the Babylonian representation
on all four models, increasing in strength roughly as
Pythia $<$ Gemma-4-E4B $<$ Llama $<$ Gemma-4-31B (by $\Delta$ helix
$R^2$).}
\end{quote}
Both measurement adjustments (wider window, extended basis) point the
same direction, which is the strongest evidence we have that the
structure is real rather than an artifact of either change.

\section{Method matters: two general lessons}
\label{sec:method-matters}

\paragraph{M1.\ The basis has bandwidth.}
Helix $R^2$ is a fitted regression onto a fixed set of periods. It is
sensitive only to structure at those periods; if a model encodes
something at a period the basis omits, helix $R^2$ will be low even
when the underlying structure is strong. The Babylonian result shows
this concretely: same data, same model, $+0.13$ helix $R^2$ from
adding one period (Table~\ref{tab:bab-extended}). The fix is to
\emph{always cross-check a basis-fit metric against a basis-free one}.
If the FFT shows a peak at a period not in the basis, helix $R^2$ is
undermeasuring the structure; add the period and refit.

\paragraph{M2.\ The window must wrap.}
To detect a period $T$ via FFT, the input range must contain several
wraps of $T$. One wrap is a single discontinuity, not a periodic
signal. The 0--99 window is fine for $T \in \{2, 5, 10\}$ ($50$, $20$,
$10$ wraps) and barely workable for $T = 100$ (one wrap, visible as a
slow trend). It is \emph{not} workable for $T = 60$: one wrap, in a
phase-misaligned position relative to the window's endpoints, produces
noise rather than a trend. The rule is to set $n_{\max}$ to a large
multiple of every candidate period.

These two corrections compose: on Babylonian at $n_{\max}=100$ with
the paper basis, both blind spots conspire in the same direction
(``no base-60 here''). Fixing only one and not the other gives
partial evidence; fixing both, and observing that the metrics move
together and in the predicted direction, is what makes the inference
robust.

\section{Discussion and conclusion}
\label{sec:discussion}

Putting the threads together:
\begin{itemize}
  \item \textbf{Cross-architecture}: the helix appears on Gemma-4-E4B
    and Gemma-4-31B---a new model family beyond the three (GPT-J,
    Pythia, Llama) the original paper tested---and re-measurement on
    Pythia and Llama under our wider protocol replicates the helix
    at very different depths and with qualitatively different
    layer-profiles. ``Read the middle layer'' is a Pythia/GPT-J
    convention. Find-the-layer is a necessary step for every new model.
  \item \textbf{Cross-script, positional}: the geometric encoding
    transfers across glyph systems for any positional base-10 script
    the model has been trained on. Multilingual models reproduce the
    helix on Devanagari and Arabic-Indic; English-heavy models do
    not. The encoding is built per-script in pre-training;
    positional structure is necessary but not sufficient.
  \item \textbf{Cross-script, non-positional}: Greek and Roman both
    collapse the helix metric, but for different reasons. Roman
    splits into a number-line module and a letter-cycling module
    living in different layers; Greek never forms a quantitative
    representation at all (the only periodic structure in the FFT
    is a tokeniser artifact at $L=0$). The two failure modes are
    distinguishable only by reading PC1 $R^2$ and its layer-of-peak
    alongside helix $R^2$. The 2-D PCA scatter
    (Section~\ref{sec:pca2d}) further reveals that Roman is encoded
    as discrete prefix-class clusters with crude ordering within each,
    and surfaces the most surprising finding of the script axis:
    on Gemma the first principal component of the residual stream
    captures up to $96.6\%$ of variance but represents
    \emph{tokenisation} (digit-count, leading-digit class), not
    magnitude. High PC1 $R^2$ is not by itself evidence of a
    magnitude axis; the helix in Gemma lives in lower-variance
    directions.
  \item \textbf{Cross-script, sexagesimal}: Babylonian cuneiform's
    base 60 is represented on all four models---a stronger
    generalisation than the paper makes, lying entirely outside the
    representational space its measurement protocol can resolve.
    Recovering it required widening the window past one wrap and
    adding $T=60$ to the trig basis.
\end{itemize}

We started this work expecting the helix to be a property of one
model's training. What appears truer now is that the helix is a
property of how transformers encode anything they represent as a
quantity, parameterised by which scripts the model has seen and which
periods are natural to those scripts. Different models pick up
different periods at different depths; the geometric strategy is the
same. Whether the helix subspace is also \emph{the representation that
downstream computation reads} is a causal question this paper does not
answer; the natural intervention sits at the answer position rather
than the operand position, and we leave it to follow-up work.

\subsection{Follow-ups}

The natural next experiments are causal:

\begin{enumerate}
  \item \textbf{Helix-subspace patching at the answer position.}
    The original paper's strongest causal claim is that on
    Pythia/Latin, patching only the helix-basis subspace
    ($\sim$9 directions of several thousand) at the position whose
    hidden state produces the answer is sufficient to control the
    answer. Replicating that on Pythia and extending to Llama,
    Gemma-4-E4B and Gemma-4-31B---and to the non-Latin positional
    scripts where the geometric helix is present in this paper---is
    the most direct causal follow-up to the work reported here.
  \item \textbf{Wider basis.} The helix basis used here is the paper's
    $T \in \{2,5,10,100\}$. Extending to the basis that recovers the
    Babylonian base-60 finding ($T \in \{2,5,10,60,100\}$), or to a
    \emph{learned} subspace whose rank is set by cumulative
    explained-variance rather than by a prescribed period list, is the
    natural follow-up if a four-period helix-subspace intervention
    leaves a large unexplained fraction of the causal effect.
  \item \textbf{Sexagesimal rotation.} On Gemma-4-31B's Babylonian
    representation at $n_{\max}=600$, rotating only the
    $(\cos 2\pi a/60, \sin 2\pi a/60)$ direction pair should shift the
    answer-token logit distribution by exactly~$60$ if the model uses
    the base-60 axis computationally. This is the cleanest test of
    representation-vs-use for a non-base-10 period.
  \item \textbf{Cross-layer units-digit rescue.} For the E4B/Latin
    dissociation: project residual at $L_{5}$ onto
    $(\cos 2\pi a/10, \sin 2\pi a/10)$ and inject the projection at a
    later layer to see whether the units-digit information is
    discarded by downstream layers or remains accessible.
\end{enumerate}

Secondary follow-ups suggested by the descriptive work but not
addressed here: Greek alphabetic numerals plausibly form $9$ clusters
of $10$, one per tens-letter, visible only via 2-D PCA scatter (the
diagnostic is included but cluster analysis is left for future work);
Roman's encoding looks like a piecewise-linear staircase with
discontinuities at glyph thresholds, which an additive regression
$h(a) \approx \sum_{\text{symbol}} \text{count(symbol)} \cdot u_{\text{symbol}}$
would test directly; and the layer at which the Babylonian base-60
representation peaks, and how it differs from the Latin base-10 peak
in the same model, is unaddressed in this paper and worth dedicated
analysis.

\section*{Reproducibility}

All code, configuration, and figures are in the public repository
accompanying this paper. Each experiment lives in its own driver
script with PEP 723 inline metadata, runnable via \texttt{uv run}; each
has an accompanying \texttt{./run\_*.sh} bash wrapper that loops over
the four models and is idempotent (each combination is skipped if its
expected output file already exists).

\paragraph{Representational sweep (Sections~\ref{sec:arch}--\ref{sec:babylonian}).}
\texttt{./run.sh} drives 40 combinations across 4 models, 8 scripts, 2
windows ($n_{\max} \in \{100, 600\}$), 2 bases ($T = \{2,5,10,100\}$
and $T = \{2,5,10,60,100\}$):
\begin{verbatim}
# Default: Pythia-6.9B, Latin digits, paper basis, n=100
uv run main.py

# Cross-script
uv run main.py --model meta-llama/Llama-3.1-8B --script devanagari
uv run main.py --model google/gemma-4-31B --script babylonian --sweep

# Wider window
uv run main.py --script babylonian --n_max 600 --sweep

# Extended basis with T = 60
uv run main.py --script babylonian --n_max 600 \
               --periods 2,5,10,60,100 --sweep
\end{verbatim}

Hardware: Apple M5 Max, 128\,GB unified memory (MPS); falls back to
CUDA, then CPU. Total wall-time for the representational sweep:
$\sim$8 hours.

\begin{thebibliography}{9}
\bibitem{kantamneni2025helix}
S.\ Kantamneni and M.\ Tegmark.
\newblock Language Models Use Trigonometry to Do Addition.
\newblock \emph{arXiv preprint} arXiv:2502.00873, 2025.

\bibitem{nanda2023progress}
N.\ Nanda, L.\ Chan, T.\ Lieberum, J.\ Smith, and J.\ Steinhardt.
\newblock Progress Measures for Grokking via Mechanistic Interpretability.
\newblock In \emph{International Conference on Learning Representations (ICLR)}, 2023.
\newblock arXiv:2301.05217.

\bibitem{power2022grokking}
A.\ Power, Y.\ Burda, H.\ Edwards, I.\ Babuschkin, and V.\ Misra.
\newblock Grokking: Generalization Beyond Overfitting on Small Algorithmic Datasets.
\newblock \emph{arXiv preprint} arXiv:2201.02177, 2022.
\end{thebibliography}

\end{document}
